%Ajustes na legenda da figura. Incluindo espacamento apos a legenda
\captionsetup[figure]{labelformat={default},labelsep=period,font=footnotesize, name=\footnotesize{Fig.},justification=centering,singlelinecheck=false,belowskip=0\normalbaselineskip}
%\pagenumbering{gobble}

%Ajustes na legenda da tabela. 
\captionsetup[table]{labelformat={default},labelsep=newline,font={sc,footnotesize},justification=centering,singlelinecheck=false}

%\renewcommand{\headrulewidth}{0pt}

% Renomear a legenda de listings para Português sem usar caption (evita conflitos)
% Para o pacote `listings` (ambiente `lstlisting`)
\renewcommand{\lstlistingname}{Código}
% Renomear também o título da lista de listings, se usado
\renewcommand{\lstlistlistingname}{Lista de Códigos}

\makeatletter
\newcommand{\linebreakand}{%
  \baselineskip 
  \end{@IEEEauthorhalign}
  \hfill\mbox{}\par
  \mbox{}\hfill\begin{@IEEEauthorhalign}
}
\makeatother


%%%%%%%%%%%%%%%%%%%%%%%%%%%%%%%%%%%%%%%%%%%%%%%%%%%%%%%%%%%%%%%%%%%%%%%%%
%    Configuracaoes de Idioma
%    Considerando babel = brazil
%%%%%%%%%%%%%%%%%%%%%%%%%%%%%%%%%%%%%%%%%%%%%%%%%%%%%%%%%%%%%%%%%%%%%%%%% 

\addto\captionsbrazil{
  \renewcommand{\abstractname}{Abstract}
  \renewcommand{\tablename}{TABELA}
}

%considerando babel = english
\addto\captionsenglish{
  \renewcommand{\tablename}{TABLE}
}